% processing.tex
\begin{python}
    from scripts.calc import *
\end{python}

\section*{\fbox{Задание 1}}
\begin{quote}
    \textit{Собрать схему и измерить все токи и напряжения ветвей. Проверить результаты, используя уравнения Кирхгофа.}
\end{quote}

Проверка первого закона Кирхгофа для узлов A и D:
\begin{equation}
    - I_1 - I_2 + I_3 
    = - \pv[-3.3]{p['ef']['I1']} - \pv[-3.3]{p['ef']['I2']} + \pv[-3.3]{p['ef']['I3']}
    = \pv[-3.3]{r['KCL_A']} \text{ мА}
\end{equation}
\begin{equation}
    I_2 + I_4 - I
    = \pv[-3.3]{p['ef']['I2']} + \pv[-3.3]{p['ef']['I4']} - \pv[-3.3]{p['ef']['I']}
    = \pv[-3.3]{r['KCL_D']} \text{ мА}
\end{equation}

Проверка второго закона Кирхгофа для левого и правого контуров (по часовой стрелке):
\begin{equation}
    U_1 + U_3 - U
    = \pv[.3]{p['ef']['U1']} + \pv[.3]{p['ef']['U3']} - \pv[.3]{p['ef']['U']}
    = \pv[.3]{r['KVL_1']} \text{ В}
\end{equation}
\begin{equation}
    -U_2 + U_4 - U_3
    = -\pv[.3]{p['ef']['U2']} + \pv[.3]{p['ef']['U4']} - \pv[.3]{p['ef']['U3']}
    = \pv[.3]{r['KVL_2_sum']} \text{ В}
\end{equation}

\section*{\fbox{Задание 2}}
\begin{quote}
    \textit{Определить методом наложения токи в ветвях и сопоставить с измеренными.}
\end{quote}

Расчет полных токов через алгебраическую сумму частичных:
\begin{equation}
    I_1
    = I_1' - I_1''
    = \pv[-3.3]{p['s']['onlyU']['I1']} - \pv[-3.3]{p['s']['onlyI']['I1']}
    = \pv[-3.3]{r['I1_sup']} \text{ мА}
\end{equation}
\begin{equation}
    I_2
    = I_2'' - I_2'
    = \pv[-3.3]{p['s']['onlyI']['I2']} - \pv[-3.3]{p['s']['onlyU']['I2']}
    = \pv[-3.3]{r['I2_sup']} \text{ мА}
\end{equation}
\begin{equation}
    I_3
    = I_3' + I_3''
    = \pv[-3.3]{p['s']['onlyU']['I3']} + \pv[-3.3]{p['s']['onlyI']['I3']}
    = \pv[-3.3]{r['I3_sup']} \text{ мА}
\end{equation}
\begin{equation}
    I_4
    = I_4' + I_4''
    = \pv[-3.3]{p['s']['onlyU']['I4']} + \pv[-3.3]{p['s']['onlyI']['I4']}
    = \pv[-3.3]{r['I4_sup']} \text{ мА}
\end{equation}

\begin{longtblr}[
    caption = {Сравнение измеренных и рассчитанных методом наложения токов},
    label = {tab:sup_comp},
    expand = \expanded,
]{
    colspec = { Q[c,m] X[c,m] X[c,m] },
    hlines, vlines,
    rowhead = 1,
}
    Ветвь & $I_{\text{изм}}$, мА & $I_{\text{расч}}$, мА \\
    \expanded{\pt*[.0, -3.3, -3.3]{
        r['sup_branches'],
        r['sup_meas'],
        r['sup_calc']
    }}
\end{longtblr}

\section*{\fbox{Задание 3}}
\begin{quote}
    \textit{Данные эксперимента подтвердить результатами расчета цепи методом эквивалентного источника.}
\end{quote}

Из методички известно, что $R_1 = \pv[.0]{r['R1']}~Ом$; $R_2 = \pv[.0]{r['R2']}~Ом$; $R_3 = \pv[.0]{r['R3']}~Ом$; $R_4 = \pv[.0]{r['R4']}~Ом$.

Найдем $R_э$:
\begin{equation}
    R_э
    = R_1 + \frac{R_3 (R_2 + R_4)}{R_3 + R_2 + R_4}
    = \pv[.0]{r['R1']} + \frac{
        \pv[.0]{r['R3']} \cdot (\pv[.0]{r['R2']} + \pv[.0]{r['R4']})
    }{
        \pv[.0]{r['R3']} + \pv[.0]{r['R2']} + \pv[.0]{r['R4']}
    }
    = \pv[.0]{r['Re']} \text{ Ом}
\end{equation}
Напряжение холостого хода $U_0$, определенное методом узловых потенциалов при отключенной ветви 3:
\begin{equation}
    U_0
    = \pv[.3]{r['U0_th']} \text{ В}
\end{equation}
Расчетный ток в ветви 3:
\begin{equation}
    I_{3\text{,расч}}
    = \frac{U_0}{R_э}
    = \frac{\pv[.3]{r['U0_th']}}{\pv[.0]{r['Re']}}
    = \pv[-3.3]{r['I3_th']} \text{ мА}
\end{equation}

Экспериментальные значения: $I_3 = \pv[-3.3]{p['mei']['I3']} \text{ мА}$. Расхождение с теорией находится в пределах погрешности.


\section*{\fbox{Задание 4}}
\begin{quote}
    \textit{Убедиться в равенстве токов при взаимном обмене источника и амперметра.}
\end{quote}

Согласно эксперименту, ток в ветви 3 при источнике в ветви 1 равен $I_3
    = \pv[-3.3]{p['pv']['I3']} \text{ мА}$. Ток в ветви 1 при переносе источника в ветвь 3 равен $I_1
    = \pv[-3.3]{p['pv']['I1']} \text{ мА}$. Равенство токов строго подтверждает принцип взаимности.

\section{Ответы на контрольные вопросы}
\subsection{Каковы результаты контроля данных в задании 1?}
Результаты контроля данных  показывают выполнение первого и второго законов Кирхгофа с точностью до тысячных долей миллиампера и вольта, что укладывается в допустимую погрешность цифровых измерительных приборов.
    
\subsection{Изменятся ли токи ветвей, если одновременно изменить полярность напряжения ИН и направление тока ИТ на противоположные?}
При одновременном изменении полярности источника напряжения и направления источника тока на противоположные, направления всех токов ветвей изменятся на инверсные, однако их абсолютные значения останутся неизменными в силу линейности цепи.

\subsection{Чему равно напряжение между узлами «C» и «D» цепи?}
Напряжение между узлами «С» и «D» определяется так:
\begin{equation}
    U_{CD}
    = U_1 - U_2
    = \pv[.3]{p['ef']['U1']} - \pv[.3]{p['ef']['U2']}
    = \pv[.3]{r['U_CD']} \text{ В}
\end{equation}

\subsection{Как изменить напряжение ИН, чтобы ток $I_1$ стал равен нулю?}
Чтобы ток $I_1$ стал равен нулю, необходимо, чтобы частичные токи от ИН и ИТ полностью компенсировали друг друга: $I_1' = I_1''$. Новое напряжение ИН рассчитывается из пропорции:
\begin{equation}
    U_{\text{нов}}
    = U \frac{I_1''}{I_1'}
    = \pv[.0]{r['U_nom']} \cdot \frac{\pv[-3.3]{p['s']['onlyI']['I1']}}{\pv[-3.3]{p['s']['onlyU']['I1']}}
    = \pv[.3]{r['U_new']} \text{ В}
\end{equation}

\subsection{Почему \cref{fig:circuit_reciprocity}б при $U=U_0$ реализует схему метода эквивалентного источника напряжения (\cref{fig:circuit_thevenin}а)?}
Схема на \cref{fig:circuit_reciprocity}б при $U = U_0$ реализует метод эквивалентного источника, так как согласно теореме Тевенена, вся активная цепь заменяется ИН, равным напряжению холостого хода, включенным последовательно с эквивалентным сопротивлением. Подключение источника $U_0$ в выделенную ветвь пассивной цепи воссоздает идентичный режим работы для данной ветви.

\subsection{Чему будет равен ток $I_1$, если ИН поместить в ветвь 4, а ИТ отключить?}
Если ИН поместить в ветвь 4, а ИТ отключить, ток $I_1$ будет равен току $I_4$ в исходной схеме при действии только ИН в силу принципа взаимности.
\begin{equation}
    I_1
    = I_4'
    = \pv[-3.3]{r['I1_Q6']} \text{ мА}
\end{equation}

\subsection{Как проконтролировать результаты экспериментов в заданиях 2--4?}
Результаты экспериментов контролируются следующими проверками: в задании 2 сумма частичных токов должна давать полные токи из задания 1; в задании 3 измеренный ток должен совпадать с током ветви 3 из задания 1; в задании 4 токи прямого и обратного измерений должны быть равны между собой.

\section*{Выводы}
В ходе лабораторной работы подтверждена справедливость законов Кирхгофа для линейных цепей. Метод наложения продемонстрировал выполнение принципа суперпозиции: алгебраическая сумма реакций от отдельных воздействий совпадает с реакцией при их совместном действии. Теорема об эквивалентном источнике Тевенена успешно применена для расчета тока в отдельной ветви, показав сходимость теоретических расчетов с результатами прямых измерений. Экспериментально доказан принцип взаимности: перестановка ИН и амперметра в пассивной цепи не изменяет значение измеряемого тока.